\begin{ex}
Số 7 gồm 4 và mấy?
\shortans{3}
\loigiai{
Ta có sơ đồ tách số: $7$ gồm $4$ và $3$ (vì $4 + 3 = 7$).
\textbf{Đáp số:} $3$.
}
\end{ex}

\begin{ex}
Số 9 gồm 2 và mấy?
\shortans{7}
\loigiai{
Ta có sơ đồ tách số: $9$ gồm $2$ và $7$ (vì $2 + 7 = 9$).
\textbf{Đáp số:} $7$.
}
\end{ex}

\begin{ex}
Số 10 gồm mấy và 6?
\shortans{4}
\loigiai{
Ta có sơ đồ tách số: $10$ gồm $4$ và $6$ (vì $4 + 6 = 10$).
\textbf{Đáp số:} $4$.
}
\end{ex}

\begin{ex}
Gộp 5 và 3 được số nào?
\shortans{8}
\loigiai{
Ta có sơ đồ gộp số: Gộp $5$ và $3$ được: $5 + 3 = 8$.
\textbf{Đáp số:} $8$.
}
\end{ex}

\begin{ex}
Gộp 4 và mấy để được 9?
\shortans{5}
\loigiai{
Để gộp 4 với một số được 9, số đó phải là: $9 - 4 = 5$.
\textbf{Đáp số:} $5$.
}
\end{ex}

\begin{ex}
Có bao nhiêu cách tách số 5 thành hai số tự nhiên khác 0 (không tính thứ tự đảo ngược)?
\shortans{2}
\loigiai{
Các cách tách số 5 thành hai số khác 0 là:\n- $5 = 1 + 4$\n- $5 = 2 + 3$\nVậy có tất cả $2$ cách tách.
\textbf{Đáp số:} $2$.
}
\end{ex}

\begin{ex}
Có bao nhiêu cách tách số 6 thành hai số tự nhiên khác 0 (không tính thứ tự đảo ngược)?
\shortans{3}
\loigiai{
Các cách tách số 6 thành hai số khác 0 là:\n- $6 = 1 + 5$\n- $6 = 2 + 4$\n- $6 = 3 + 3$\nVậy có tất cả $3$ cách tách.
\textbf{Đáp số:} $3$.
}
\end{ex}

\begin{ex}
Số 8 được tách thành hai số giống nhau. Hỏi mỗi số đó là số mấy?
\shortans{4}
\loigiai{
Vì hai số giống nhau có tổng bằng 8, mà $4 + 4 = 8$ nên mỗi số đó là số $4$.
\textbf{Đáp số:} $4$.
}
\end{ex}

\begin{ex}
Số 10 được tách thành hai số giống nhau. Hỏi mỗi số đó là số mấy?
\shortans{5}
\loigiai{
Vì hai số giống nhau có tổng bằng 10, mà $5 + 5 = 10$ nên mỗi số đó là số $5$.
\textbf{Đáp số:} $5$.
}
\end{ex}

\begin{ex}
Bạn Lan có 7 cái kẹo, Lan chia vào hai túi sao cho túi thứ nhất có 3 cái kẹo. Hỏi túi thứ hai có bao nhiêu cái kẹo?
\shortans{4}
\loigiai{
Số kẹo ở túi thứ hai là: $7 - 3 = 4$ (cái kẹo).
\textbf{Đáp số:} $4$.
}
\end{ex}

\begin{ex}
Một cây có 9 quả táo. Gió thổi làm rụng 4 quả xuống đất. Hỏi trên cây còn lại bao nhiêu quả táo?
\shortans{5}
\loigiai{
Tách 9 quả táo thành 4 quả rụng và số quả còn lại trên cây: $9 - 4 = 5$ (quả táo).
\textbf{Đáp số:} $5$.
}
\end{ex}

\begin{ex}
Gộp 3 chiếc bút chì đỏ và 6 chiếc bút chì xanh được tất cả bao nhiêu chiếc bút chì?
\shortans{9}
\loigiai{
Gộp số bút chì đỏ và số bút chì xanh: $3 + 6 = 9$ (chiếc bút chì).
\textbf{Đáp số:} $9$.
}
\end{ex}

\begin{ex}
Số 6 gồm 1 và một số. Số đó là số mấy?
\shortans{5}
\loigiai{
Ta có: $6$ gồm $1$ và $5$ (vì $1 + 5 = 6$).
\textbf{Đáp số:} $5$.
}
\end{ex}

\begin{ex}
Số 8 gồm 5 và mấy?
\shortans{3}
\loigiai{
Ta có: $8$ gồm $5$ và $3$ (vì $5 + 3 = 8$).
\textbf{Đáp số:} $3$.
}
\end{ex}

\begin{ex}
Có 10 quả cam chia đều vào hai đĩa. Hỏi mỗi đĩa có bao nhiêu quả cam?
\shortans{5}
\loigiai{
Chia đều 10 quả cam vào 2 đĩa: mỗi đĩa có $5$ quả cam (vì $5 + 5 = 10$).
\textbf{Đáp số:} $5$.
}
\end{ex}

\begin{ex}
Một nhóm gồm 8 bạn học sinh được tách thành hai nhóm nhỏ: nhóm nam và nhóm nữ. Biết nhóm nữ có 3 bạn, hỏi nhóm nam có bao nhiêu bạn?
\shortans{5}
\loigiai{
Số bạn học sinh nam là: $8 - 3 = 5$ (bạn).
\textbf{Đáp số:} $5$.
}
\end{ex}

\begin{ex}
Gộp 2 que tính và 7 que tính được tất cả bao nhiêu que tính?
\shortans{9}
\loigiai{
Gộp 2 và 7 que tính: $2 + 7 = 9$ (que tính).
\textbf{Đáp số:} $9$.
}
\end{ex}

\begin{ex}
Số 9 gồm hai số tự nhiên liên tiếp. Hỏi số lớn hơn trong hai số đó là số mấy?
\shortans{5}
\loigiai{
Hai số tự nhiên liên tiếp có tổng bằng 9 là $4$ và $5$ (vì $4 + 5 = 9$ và $5 - 4 = 1$).\nSố lớn hơn trong hai số đó là số $5$.
\textbf{Đáp số:} $5$.
}
\end{ex}

\begin{ex}
Số 7 gồm hai số tự nhiên liên tiếp. Hỏi số bé hơn trong hai số đó là số mấy?
\shortans{3}
\loigiai{
Hai số tự nhiên liên tiếp có tổng bằng 7 là $3$ và $4$ (vì $3 + 4 = 7$).\nSố bé hơn trong hai số đó là số $3$.
\textbf{Đáp số:} $3$.
}
\end{ex}

\begin{ex}
Gộp 1, 3 và 5 được số nào?
\shortans{9}
\loigiai{
Gộp ba số lại: $1 + 3 + 5 = 9$.
\textbf{Đáp số:} $9$.
}
\end{ex}

\begin{ex}
Số 10 gồm 3 số: 2, 5 và một số nữa. Hỏi số còn lại là số mấy?
\shortans{3}
\loigiai{
Tổng của hai số đã biết là: $2 + 5 = 7$.\nSố còn lại là: $10 - 7 = 3$.
\textbf{Đáp số:} $3$.
}
\end{ex}

\begin{ex}
Có 8 bông hoa cắm vào hai lọ sao cho số hoa ở hai lọ bằng nhau. Hỏi mỗi lọ có bao nhiêu bông hoa?
\shortans{4}
\loigiai{
Vì số hoa ở hai lọ bằng nhau và $4 + 4 = 8$ nên mỗi lọ có $4$ bông hoa.
\textbf{Đáp số:} $4$.
}
\end{ex}

\begin{ex}
Tách số 7 thành hai số có hiệu bằng 3. Hỏi số lớn hơn là số mấy?
\shortans{5}
\loigiai{
Hai số có tổng bằng 7 và hiệu bằng 3 là $5$ và $2$ (vì $5 + 2 = 7$ và $5 - 2 = 3$).\nSố lớn hơn là số $5$.
\textbf{Đáp số:} $5$.
}
\end{ex}

\begin{ex}
Tách số 6 thành hai số có hiệu bằng 4. Hỏi số bé hơn là số mấy?
\shortans{1}
\loigiai{
Hai số có tổng bằng 6 và hiệu bằng 4 là $5$ và $1$ (vì $5 + 1 = 6$ và $5 - 1 = 4$).\nSố bé hơn là số $1$.
\textbf{Đáp số:} $1$.
}
\end{ex}

\begin{ex}
Gộp 4 bạn nam và 6 bạn nữ lại thành một đội văn nghệ. Hỏi đội văn nghệ đó có tất cả bao nhiêu bạn?
\shortans{10}
\loigiai{
Gộp số bạn nam và nữ: $4 + 6 = 10$ (bạn).
\textbf{Đáp số:} $10$.
}
\end{ex}
